\subsection*{A completed comparison across one actual neck}

\statement{B.108}{Proved}{The actual first-neck slit model has a proper full profinite amalgam and a vertex centralizer reduction}
\label{stmt:B108}
Use the actual surface $\mathscr Q$, involution $F(q,t)=(q,-t)$,
base point $\bar P$ and projective functions $T_u$ of B.106.
Write $u_j=\exp(2\pi i j/5)$, $1\leq j\leq4$, and fix
$q_*=\exp(\pi i/5)$. Choose a sufficiently small disk $N$ about
$q_*$ avoiding all old punctures and all critical values. Choose
simply connected slit regions $V_-$ and $V_+$ in $0<|q|<1$ and
$1<|q|<\infty$, respectively, retaining all their nine critical
values, with $\bar P$ over $V_-$. The slits and the enlargements
through $N$ are fixed once and for all, and avoid the indicated
finite sets. Their union with this one neck is denoted by $U_*$;
the two enlarged regions have one contractible overlap. No other
unit-circle neck or old base loop is included.

Let $R_-$ be the orbifold fundamental group of the bad region
$|t|>1$ over $V_-$, and $R_+$ that of $|t|<1$ over $V_+$,
using the variable collars of B.107. An orbifold fundamental group
here means the fundamental group of the Borel construction for
the specified $C_2$ action. Choose all basepoint connectors through
the same regular neck. Its disk subgroup
\[
 K=\pi_1(D^-\setminus\{T_{u_1},\ldots,T_{u_4},q_*^{-1/2}\})
       \cong F_5
\]
is specified in the proof: $D^-$ lies in the lower $t$ semicircle,
and $q_*^{-1/2}=\exp(-\pi i/10)$. Then the actual based orbifold
group and its full profinite completion satisfy
\begin{equation}\label{eq:B108-amalgams}
 \begin{aligned}
 G_*:=\pi_1(\mathscr Q_{U_*},\bar P)\rtimes_F C_2
      &\cong R_-*_{K}R_+,\\
 \widehat G_*&\cong
      \widehat R_-\amalg_{\widehat K}\widehat R_+.
 \end{aligned}
\end{equation}
The discrete edge maps and both maps of vertex completions in the
second line are injective. The second line is a proper amalgam
in the category of all profinite groups.

For either vertex, let $H_\pm$ be its ordinary index-two subgroup
and $\tau_\pm$ its actual order-two cone. There are homomorphisms
\begin{equation}\label{eq:B108-paired-retractions}
 r_\pm:H_\pm\longrightarrow K,
 \qquad r_\pm|_K=\operatorname{id},\qquad
 r_\pm(\tau_\pm K\tau_\pm)=1.
\end{equation}
In particular every finite-index normal subgroup of $K$ is the
exact intersection with $K$ of a finite-index normal subgroup of
each vertex. The cone centralizers obey
\begin{equation}\label{eq:B108-centralizers}
 C_{G_*}(\tau_\pm)=C_{R_\pm}(\tau_\pm),\qquad
 C_{\widehat G_*}(\tau_\pm)
       =C_{\widehat R_\pm}(\tau_\pm).
\end{equation}
The selected cone at $t=0$ belongs to $R_+$ after the neck
identification. Its centralizer in that vertex is not evaluated
by this statement. No full-surface fixed-source equality follows.

\begin{proof}
\emph{B.108.1. One actual common marking.}
Put
\[
 Q=\frac{1+q}{1-q},\qquad Z=\frac{t+1}{t-1},\qquad
 c_j=\cot(\pi j/5),\qquad \lambda_j=1-\cos(2\pi j/5)>0.
\]
Substitution in B.106 gives
\[
 Z(T_{u_j})=\lambda_j(Q-i c_j).
\]
At $q=\exp(i\theta)$ with $0<\theta<2\pi/5$, this equals
$i\lambda_j(\cot(\theta/2)-c_j)$ and has positive imaginary
part for every $j$. Since
\[
 \operatorname{Im}\frac{Z+1}{Z-1}
       =\frac{-2\operatorname{Im}Z}{|Z-1|^2},
\]
all four literal $A$ punctures $T_{u_j}$ lie in the lower $t$
semicircle. The branch puncture $\exp(-i\theta/2)$ lies there
also; their five $F$ partners lie in the upper semicircle.
All are distinct, by the collision factorizations in B.106.

In the orbifold quotient coordinate $y=t^2$, the fiber has cones
$0,\infty$ and the six ordinary punctures
$1,1/q,T_{u_1}^2,\ldots,T_{u_4}^2$. At $q_*$ choose an
annular-sector disk $D$ containing the five punctures other than
$1$, and avoiding both cones and $1$. Its two lifts are disks;
the lower lift $D^-$ contains exactly the five punctures in the
statement. Fix disjoint meridian arcs there and one connector
through the gap adjoining the retained puncture. These choices
specify the literal based $K$, rather than only its conjugacy class.

Over $N$ the actual puncture motion transports this disk and its
arc system. At a nearby inner endpoint all moving punctures have
$|t|>1$; at a nearby outer endpoint they have $|t|<1$. These
inequalities are B.107's exact separation and its outer mirror.
At each endpoint the transported lower-half-plane disk can be
isotoped into the corresponding bad side, with its connector in
the fixed gap. To see this without a condition at the middle
fiber, work in the lower half-plane, which contains exactly the
five selected punctures. Both endpoint representatives bound a
disk enclosing all five; an isotopy of the punctured half-plane,
fixed at the gap, identifies them. No disk is required to be in
an open bad side at $|q|=1$ itself.

The common orbifold fiber marking is therefore
\[
 C_{2,0}*C_{2,\infty}*K.
\]
Over the simply connected slit regions, the good-side factors
of B.107 are just their cones. Orbifold van Kampen gives
$E_-=C_{2,0}*R_-$ and $E_+=C_{2,\infty}*R_+$.
In their pushout the two free cone generators are identified
with the actual opposite bad-side cones. Eliminating those two
redundant generators leaves the pushout $R_-*_{K}R_+$.
This elimination does not assume injection of the full neck
fiber group. We prove the required $K$ injections next.

\emph{B.108.2. The actual paired retractions.}
Fill only the four $B$ divisors $t=-T_u(q)$ in the ordinary
surface. The result is the actual bundle $\mathscr Q_A$ with
fiber a sphere punctured at
\[
 1,-1,\ q^{-1/2},-q^{-1/2},\ T_{u_1},\ldots,T_{u_4}.
\]
All possible coincidences of these remaining punctures lie over
the already deleted old base values, so this bundle includes
every critical value retained in either slit region. Its fiber
group is $F_7$. Because the slit base has both $\pi_1=0$ and
$\pi_2=0$, the based fiber inclusion induces an isomorphism on
fundamental groups. Composing ordinary inclusion and filling
therefore gives a homomorphism $f_\pm:H_\pm\to F_7$.
The construction is on the ordinary index-two subgroup only;
filling one family is not $F$-equivariant.

The chosen disk contains five of these eight punctures, and
the other branch and the two retained punctures lie outside.
Extend its already chosen geometric meridian basis by a meridian
$b$ of the other branch and a meridian $d$ of one retained point;
omit the remaining sphere meridian. Disjoint outside arcs give
\[
 F_7=K*F(b,d).
\]
The projection $p:F_7\to K$ fixing $K$ and killing $b,d$ gives
$r_\pm=p f_\pm$, with restriction the identity on $K$.
In particular the actual edge maps into $H_\pm$, and hence into
$R_\pm$, are injective. This also justifies the ordinary amalgam
interpretation of the pushout above.

Let $\alpha_\pm$ be conjugation by $\tau_\pm$ on $H_\pm$.
Each of the four $A$ generators of $K$ is taken to a based
conjugate of a $B$ peripheral. Filling kills its whole normal
closure. The fifth generator is taken to a based conjugate of
the other branch peripheral, whose whole normal closure is
killed by $p$. All whiskers are retained in these assertions:
the image of a conjugate $h m h^{-1}$ is trivial because the
image of its core $m$ is trivial, not because $h$ is discarded.
A common change from the global connector to the vertex cone
only changes those conjugators. Thus
$r_\pm\alpha_\pm(K)=1$, proving
\eqref{eq:B108-paired-retractions}.

The ordinary homomorphism
\[
 h\longmapsto(r_\pm(h),r_\pm(\alpha_\pm(h)))
\]
is equivariant for interchange of the two factors. It extends to
\begin{equation}\label{eq:B108-paired-map}
 R_\pm\longrightarrow (K\times K)\rtimes_{\mathrm{swap}}C_2,
 \qquad k\longmapsto(k,1),\quad
 \tau_\pm\longmapsto\text{the interchange involution}.
\end{equation}
For any finite-index normal $L\triangleleft K$, compose with
$K/L$ in both factors. The resulting finite quotient retains
orientation and has kernel intersecting $K$ in exactly $L$.

\emph{B.108.3. Compatible finite quotients and proper completion.}
Given prescribed finite-index normal subgroups
$M_\pm\triangleleft R_\pm$, choose
$L\triangleleft K$ of finite index contained in both
$M_\pm\cap K$. Intersect $M_\pm$ with the respective kernels
of \eqref{eq:B108-paired-map} modulo $L$, obtaining
$N_\pm\triangleleft R_\pm$. Then
\[
 N_\pm\subset H_\pm,\qquad N_\pm\cap K=L.
\]
The finite groups $A_\pm=R_\pm/N_\pm$ have the same embedded
edge $D=K/L$. Form the discrete finite-group amalgam
\[
 V=A_-*_{D}A_+.
\]
We record explicitly the elementary finiteness facts required
below. Let a finite set have cardinality a common multiple of
$|A_-|$ and $|A_+|$. Realize each vertex by copies of its regular
action. Their restrictions are isomorphic free $D$-sets; identify
them. The resulting homomorphism from $V$ to the finite symmetric
group is injective on both vertices. Its kernel acts freely on
the ordinary amalgam tree and is free. It has finite rank because
it has finite index in the finitely generated group $V$.
Residual finiteness of that free kernel, followed by finite-index
cores in $V$, proves residual finiteness of $V$ itself.

Every finite quotient of $G_*$ factors through some such $V$:
apply the construction to the kernels of its two vertex
restrictions. Conversely, the finite permutation quotient of
$V$ retains both finite vertex images. Thus every prescribed
finite quotient of either $R_\pm$ is retained after refinement
in a finite quotient of $G_*$. Both vertices, and the common
edge, have their full induced profinite topologies. This proves
injectivity of their completed maps. The universal property on
finite targets gives the second line of
\eqref{eq:B108-amalgams}, with properness now proved. We have
not assumed residual finiteness of the unknown $R_\pm$.

\emph{B.108.4. The prime-order theorem has all its hypotheses.}
For each finite amalgam $V$, fix a normal finite-index free
subgroup $F_r$ with finite quotient $D_0$. Finite-index cores
give the full topology on $F_r$, so completion yields the exact
extension with kernel $\widehat F_r$ and quotient $D_0$.
For every discrete continuous
$\mathbb F_p[[\widehat F_r]]$-module $A$, the cohomology
comparison is an isomorphism in all degrees. Degree zero
follows from density. Degree one follows by prescribing a
cocycle arbitrarily on the finite free basis: its values and
their compact-group orbits span a finite invariant submodule,
and the corresponding finite semidirect product gives the
continuous extension of the cocycle. Coboundaries agree.
Higher degrees vanish because both free groups have
$p$-cohomological dimension at most one.

The Lyndon--Hochschild--Serre comparison for the two extensions
over $D_0$ has identical $E_2$ pages and only rows zero and
one. It proves the all-degree comparison for every discrete
continuous $\mathbb F_p[[\widehat V]]$-module. With finite
coefficients the two free-group rows are finite, as are their
cohomology groups under the finite group $D_0$. The same
spectral sequence makes $H^n(V,A)$ finite in every degree.
Also $\operatorname{vcd}_p(V)\leq1$. Thus $V$ is residually
finite, good for the required arbitrary discrete modules, and
of finite virtual $p$-cohomological \emph{type}, including its
finiteness condition. In particular every hypothesis of
\citet[Theorem~4.3, prime-order case in Sections~4.1--4.2]{BoggiZalesskii2025}
holds at $p=2$. We use only the element-centralizer conclusion
for the order-two cone, whose complete proof was reopened at
this dependency.

\emph{B.108.5. Centralizers and the cofinal separation argument.}
In the ordinary amalgam $G_*$, an involution in a vertex cannot
fix an edge because $K$ is torsion-free. Its fixed subtree is
therefore exactly that vertex, giving the discrete equality
in \eqref{eq:B108-centralizers}. In a finite amalgam $V$, the
image of $\tau_\pm$ still has order two and odd orientation,
whereas every conjugate of the finite edge has even orientation.
It again fixes precisely its vertex. Hence
\[
 C_V(\tau_\pm)=C_{A_\pm}(\tau_\pm).
\]
The right side is finite. The verified prime-order theorem and
the hypotheses just proved give
\[
 C_{\widehat V}(\tau_\pm)=C_{A_\pm}(\tau_\pm).
\]

Let $g\in C_{\widehat G_*}(\tau_\pm)$. Its image in every
$\widehat V$ belongs to the corresponding $A_\pm$. If $g$
lay outside the compact image of $\widehat R_\pm$, some finite
quotient of $\widehat G_*$ would separate it from that image.
That quotient factors through one of the $V$ just constructed,
contradicting the stated membership in its finite vertex. Thus
$g\in\widehat R_\pm$, and the completed equalities follow.
This argument makes no choice of compatible centralizer
representatives and completes no arbitrary infinite-index
kernel. The actual remaining vertex centralizer, the other
four necks, and the old base loops remain separate questions.
\end{proof}
